\documentclass[11pt]{article}
\usepackage[margin=1in]{geometry}
\usepackage{amsmath,amssymb}
\usepackage[hidelinks]{hyperref}

\newcommand{\smin}{\sigma_{\min}}

\title{A later-literature answer to the $k=2$ subcase\\of Conjecture 5.2 in arXiv:1308.4718}
\author{Literature-status packet for run \texttt{fa\_banach\_001}}
\date{27 August 2026}

\begin{document}
\maketitle

\paragraph{Status.}
\texttt{literature\_implied\_answer\_partial\_subcase}.  This note records a
known later result, not a new solution produced by the run.

\section*{Source question}

Balan--Wang consider, on PDF page 16 of arXiv:1308.4718v1, an
$n\times(n+k)$ real matrix
\[
 F=[g_1,\ldots,g_{n+k}], \qquad \|g_j\|\le 1,
\]
and define
\[
 \tau(F)=\min_{\substack{S\subseteq\{1,\ldots,n+k\}\\|S|=n}}
 \sigma_n(F_S).
\]
Their Conjecture 5.2 asks whether, for fixed $k$, there is a constant
$C=C(k)$ such that
\begin{equation}\label{eq:source}
 \tau(F)\le \frac{C}{n^{k-1/2}}.
\end{equation}

\section*{Later theorem and exact identification}

Liu--Wang prove in Lemma 2.4 of their 2016 paper that every real
$(n+2)\times n$ matrix $A$ whose rows have norm at most one satisfies
\begin{equation}\label{eq:later}
 \min_{\substack{S\subseteq\{1,\ldots,n+2\}\\|S|=n}}
 \sigma_n(A_S)\le \frac{C}{n^{3/2}}
\end{equation}
for an absolute constant $C$.

Take $k=2$ in the source and set $A=F^T$.  The rows of $A$ are precisely
the columns $g_j$ of $F$, so they have norm at most one.  For every
$n$-element index set $S$,
\[
 A_S=F_S^T
 \quad\text{and hence}\quad
 \sigma_n(A_S)=\sigma_n(F_S).
\]
Thus \eqref{eq:later} is exactly \eqref{eq:source} for $k=2$.

The later proof uses a two-dimensional orthonormal basis of the kernel
and reduces the estimate to finding two nearby projective directions on
the circle.  This is the same direct kernel-duality route suggested by
the source problem.  Liu's 2024 paper cites both the source and the 2016
paper and gives a sharp version of the planar two-column extremal bound.

\section*{Scope and provenance}

This is a complete answer only for the $k=2$ subcase.  It does not settle
Conjecture 5.2 for arbitrary fixed $k$ and says nothing about Conjecture
5.1.  The 2016 paper's higher-dimensional discussion reduces the next
case to Heilbronn-type incidence geometry and gives partial results under
additional geometric hypotheses.

The relation is classified as a literature-implied answer rather than an
explicitly announced solution: the later paper proves the identical
matrix inequality but does not refer to it by the source's conjecture
label.  Yang Wang is a coauthor of both papers.

\begin{thebibliography}{9}
\bibitem{BW}
R. Balan and Y. Wang,
\emph{Invertibility and Robustness of Phaseless Reconstruction},
arXiv:1308.4718v1, Section 5, Conjecture 5.2, PDF page 16;
later published in Appl. Comput. Harmon. Anal. 38 (2015), 469--488.

\bibitem{LW}
Y. Liu and Y. Wang,
\emph{On the Decay of the Smallest Singular Value of Submatrices of Rectangular Matrices},
Asian-European J. Math. 9 (2016), article 1650075,
doi:10.1142/S1793557116500753, Lemma 2.4.

\bibitem{Liu2024}
Y. Liu,
\emph{On Minimal Smallest Singular Value of Subframes for Signal Recovery},
Operators and Matrices 18 (2024), 273--293,
doi:10.7153/oam-2024-18-17, Theorem 2.7.
\end{thebibliography}

\end{document}

